% Source-derived chapter 01: Introduction
% Original source: virasoro-constraints-moduli-sheaves-vertex-algebras
\chapter{Introduction}
\label{virasoro-constraints-moduli-sheaves-vertex-algebras-chapter-01}
\source{virasoro-constraints-moduli-sheaves-vertex-algebras}

\begin{remark}[Source section]
\label{virasoro-constraints-moduli-sheaves-vertex-algebras-chapter-01-source}
\source{virasoro-constraints-moduli-sheaves-vertex-algebras}
\source{virasoro-constraints-moduli-sheaves-vertex-algebras}
This chapter reproduces the corresponding section of the retrieved
LaTeX source, including its definitions, statements, examples, and
proofs. It has not yet been translated into Lean.
\notready
\end{remark}

% The source section title was: Introduction
This paper concerns the Virasoro constraints on sheaf counting theories. Given a moduli space of sheaves $M$ with a virtual fundamental class $[M]^\vir$ we may produce numerical invariants by integrating natural cohomology classes -- called descendents -- against the virtual fundamental class. The Virasoro conjecture predicts that these numerical invariants are constrained by some explicit and universal relations. The main result of this paper is proof of those constraints in the following cases: 

\begin{thm}[=Theorem \ref{thm: main result of the paper}]
\source{virasoro-constraints-moduli-sheaves-vertex-algebras}
\notready\label{thm: main} 
\source{virasoro-constraints-moduli-sheaves-vertex-algebras}
The Virasoro constraints (Conjecture \ref{conj: virasorosheaves}) hold for the following moduli spaces when there are no strictly semistable sheaves:
\begin{enumerate}
\item The moduli spaces $M_C(r,d)$ of slope semistable bundles on any smooth projective curve $C$.
\item The moduli spaces $M_S^H(r,c_1, c_2)$ of slope or Gieseker semistable torsion-free sheaves on any smooth projective surface $S$ with $h^{1,0}(S)=h^{2,0}(S)=0$.
\item  Assuming that a necessary wall-crossing formula holds (see Assumption \ref{ass:WCpair}), the moduli spaces $M_S^H(\beta, n)$ of slope semistable one-dimensional sheaves on any smooth projective surface $S$ with $h^{1,0}(S)=h^{2,0}(S)=0$.\footnote{In \cite{Pi22}, the ring structure of the cohomology of the moduli spaces of one-dimensional sheaves
on $S = \BP^2$ is studied in terms of descendents. It would be interesting to find out whether the
Virasoro constraints are implied by the relations they prove in \cite[Proposition 2.6]{Pi22}.}
\end{enumerate}

\end{thm}
%All three of these results are consequences of Theorem \ref{thm: rankreduction} which states that each of them is equivalent to checking Virasoro constraints for moduli spaces of stable pairs generalizing Pandharipande--Thomas stability. In Section \ref{sec: limits}, we discuss that many of these pair moduli spaces are trivial which turns Theorem \ref{thm: rankreduction} into a procedure that reduces rank. Virasoro constraints for low-rank moduli spaces are either known to be satisfied or we prove them in Section \ref{sec:lowrank}.

Theorem \ref{thm: main result of the paper} also proves an analog of (1), (2) and (3) in the presence of strictly semistable sheaves as stated in Conjecture \ref{conj:actualinvvir}. The case (2) of the theorem solves the conjecture of D. van Bree in \cite{bree} (see Remark \ref{rem: normalizedpt} for a comparison of our formulation with van Bree's). For $S=\PP^2$ and $\PP^1\times\PP^1$, an alternative proof of (2) with Gieseker stability was given by the first author in \cite{BoVir} without relying on the previous result for $\Hilb^n(S)$ in \cite{moop}. The formulation of the other two cases is new; indeed, we provide a very general conjecture that includes many other interesting cases. The proof of Assumption \ref{ass:WCpair} making (3) self-contained will be addressed in the future. 






Our work relies in a fundamental way on the vertex algebra that D. Joyce recently introduced \cite{Jo17,GJT, Jo21} to study the wall-crossing of moduli spaces of sheaves. We explain how the Virasoro constraints can be naturally formulated in this language.

\begin{thm}[=Theorem \ref{thm: duality} and Theorem \ref{thm: virasorophysical}]
\source{virasoro-constraints-moduli-sheaves-vertex-algebras}
\notready\label{thm: B}
\source{virasoro-constraints-moduli-sheaves-vertex-algebras}
Let $X$ be a curve or a surface with $p_g=0$. There is a natural conformal element $\omega$ in the vertex algebra $V_\bullet^{\pa}$ for which the corresponding Virasoro operators $L_n^{\pa}$ are dual to the Virasoro operators $\bL_n^{\pa}\colon \BD^{X,\pa}\to \BD^{X,\pa}$ for $n\geq -1$ on the pair descendent algebra. A moduli space of sheaves (or pairs) satisfies the Virasoro constraints if and only if its class in $\widecheck V_\bullet^{\pa}$ (or $V_\bullet^{\pa}$) is a primary state.
\end{thm}


As a consequence of this description, the wall-crossing machinery developed by Joyce can be used to prove a compatibility between wall-crossing and the Virasoro constraints.
\begin{thm}[=Propositions \ref{prop: wallcrossingcompatibility1} and \ref{prop: wallcrossingcompatibility2}]
\source{virasoro-constraints-moduli-sheaves-vertex-algebras}
\notready\label{thm: C}
\source{virasoro-constraints-moduli-sheaves-vertex-algebras}
The Virasoro constraints are compatible with wall-crossing.
\end{thm}
Theorem \ref{thm: C} also holds in the setting of quivers as was shown by the first author in \cite[Theorem 1.2]{BoVir}. This was used to conclude Virasoro constraints for moduli spaces of semistable (framed) representations of quivers in loc. cit.

Our proof of Theorem \ref{thm: main} uses this wall-crossing compatibility to reduce the statement to the rank 1 case; for curves, the rank 1 case can be shown directly or via Joyce-Song wall-crossing, and for surfaces, it was already proven in \cite{moop, moreira}. We remark that the hypotheses $h^{1,0}(S)=h^{2,0}(S)=0$ are currently necessary: $h^{2,0}(S)=0$ is needed to have a vertex algebraic interpretation of the Virasoro constraints (cf. Theorem \ref{thm: B}), and hence wall-crossing compatibility, and $h^{1,0}(S)=0$ is necessary in the proof of the constraints for the Hilbert scheme of points in \cite{moreira}; we hope these restrictions can be lifted in the future.


\subsection{History}
\label{sec: history of Virasoro constraints }\hfill
\source{virasoro-constraints-moduli-sheaves-vertex-algebras}

\noindent\textbf{Virasoro constraints} \
The study of Virasoro constraints on curve counts traces back to the origins of Gromov-Witten (GW) theory and intersection theory on the moduli space of stable curves, in Witten's foundational paper \cite{witten}. Witten conjectured that integrals of products of descendents -- certain natural classes in $H^\bullet(\overline \M_{g,n})$ -- obeyed some explicit relations. The relations he proposed were equivalent to certain differential operators,  which satisfy the Virasoro bracket relation, annihilating the partition function which encodes all the integrals of descendents on $\overline \M_{g,n}$. Witten's conjecture was proven by Kontsevich \cite{kontsevich} and new proofs were obtained by Okounkov-Pandharipande \cite{OPwitten} and Mirzakhani \cite{mirzakhani}.

In \cite{ehx}, the authors extended the Virasoro conjecture to the GW invariants of a variety $X$. Since then, a lot of effort has been put into proving the result for some target varieties $X$; most notably, the conjecture is now known when $X$ is a toric variety (or, more generally, when $X$ has semisimple quantum cohomology) by work of Givental and Teleman \cite{givental, teleman} and when $X$ is a curve by work of Okounkov-Pandharipande \cite{OP}. The general case, however, is still out of reach. 

In \cite{mnop1, mnop2}, Maulik-Nekrasov-Okounkov-Pandharipande propose a \allowbreak deep connection between Gromov-Witten invariants and Donaldson-Thomas (DT) invariants of 3-folds. Such correspondence suggested that the DT descendent invariants should as well be constrained by some sort of Virasoro operators. Almost 15 years ago, not long after the proposal of the MNOP conjecture, Oblomkov, Okounkov and Pandharipande were able to predict the precise form for the DT Virasoro operators (at least for $X=\PP^3$, see \cite[Conjecture 8]{survey}) from experimental data with $X$ toric. The understanding of the MNOP correspondence at the time, however, was not sufficiently explicit to be used effectively to relate the conjectures on the GW and on the DT sides. Recently, the GW/PT descendent correspondence has been made more effective and this allowed a proof of the DT Virasoro conjecture when $X$ is a toric 3-fold in the stationary regime \cite{moop}.\footnote{The results in \cite{moop} are formulated entirely in the theory of stable pairs, also known as Pandharipande-Thomas (PT) invariants, and not DT invariants. In the stationary regime for toric 3-folds, however, the two formulations of Virasoro are known to be equivalent. } 

Taking a surface $S$ and $X=S\times \PP^1$ it is possible to deduce some Virasoro constraints for the Hilbert scheme of points on $S$ from the PT Virasoro constraints on $X$. The third author used a universality argument in \cite{moreira} to prove such constraints for every surface with $H^1(S)=0$ by starting with the toric results in \cite{moop}. Subsequently, van Bree proposed a generalization of the Hilbert scheme constraints to the moduli spaces of torsion-free stable sheaves on a surface $S$ and made several non-trivial checks for toric $S$ using localization \cite{bree}.


While the Virasoro constraints on sheaf-counting theories come historically from Gromov-Witten theory, they form a rich theory themselves as indicated by the examples where they can be studied -- DT, PT, Hilbert scheme, stable torsion-free sheaves on surfaces. We show that they have an independent meaning and origin by connecting them to the geometric construction of the vertex algebras of Joyce \cite{Jo17} which were developed to study wall-crossing.



\noindent\textbf{Wall-crossing and vertex algebras} \
When Donaldson \cite{Do83} introduced his invariants counting anti-self-dual instantons, they were intrinsically dependent on the choice of a metric $g$ of the underlying four-manifold. Varying $g$ leads to discontinuous jumps of the invariants along codimension one walls, a phenomenon called ``wall-crossing". The precise description of the wall-crossing contributions has been given in \cite{KoMo94} and many further studies have been conducted.



With the goal of treating wall-crossing phenomena uniformly, Joyce \cite{JO06I, JO06II, JO06III, JO06IV} developed a theory which could be applied in large generality to abelian categories. Here the metric was replaced by stability conditions and instantons by semistable objects. Using a Lie algebra structure, he was able to define motivic invariants counting semistable objects and described how they change when varying the stability conditions. Further refinements to include DT theory of Calabi-Yau 3-folds were considered by Joyce-Song \cite{JS12} and Kontsevich-Soibelman \cite{KS08}.


The (virtual) fundamental classes of sheaves are however not motivic outside of the realm of Calabi-Yau 3-folds, so their theories were not sufficient for studying other geometries. In a more recent development, Joyce \cite{Jo17} introduced a sheaf-theoretic construction of vertex algebras (see \cite{Borcherds, Ka98, LLVA} for a gentle introduction to this topic).\footnote{See also H. Liu's work \cite{Liu} where he extends it to $K$-theory and multiplicative vertex algebras.} Vertex algebras are representation theoretic objects introduced by Borcherds \cite{Borcherds} and they give an axiomatization of conformal field theories in two dimensions \cite{BPZ}. The Lie bracket operation induced from the sheaf-theoretic vertex algebras was used to describe wall-crossing of virtual fundamental classes counting semistable objects, as conjectured in \cite{GJT} and proven in many cases by Joyce \cite{Jo21}. For surfaces, these wall-crossing formulae are related to the work of Mochizuki \cite{mochizuki} where the formulae are presented without vertex algebras. 

Wall-crossing has been used in \cite{Bo21.5, Bo21} to give explicit formulae for all descendent invariants of punctual Quot schemes on surfaces and Calabi-Yau fourfolds, and in \cite{Bu22} to study moduli spaces of vector bundles on curves. However, further structures coming from the vertex algebra remained a mystery. We fill this gap by giving a geometric interpretation of a natural conformal element in terms of Virasoro constraints. The conformal element  induces a representation of the Virasoro algebra on Joyce's vertex algebra \cite{Jo17}. The Virasoro operators $\{L_n\}_{n\in \BZ}$ act on the homology of the stack where wall-crossing takes place, defining a smaller Lie algebra of \textit{primary/physical states} \cite{BPZ,Borcherds}. We show that the Virasoro constraints are precisely the statement that (virtual) fundamental classes of moduli of semistable sheaves are primary states, and thus are preserved by wall-crossing. We use this new technique, together with a rank reduction to the case of rank 1, to prove existing and new conjectures about Virasoro constraints.


\subsection{Moduli of sheaves and pairs}
\label{sec: modulisheavespairs}
\source{virasoro-constraints-moduli-sheaves-vertex-algebras}

Let $X$ be a smooth projective variety over the complex numbers. Typically, we will restrict ourselves to small dimension $X$ (up to dimension 3) so that the moduli spaces of sheaves that we consider have a virtual fundamental class in the sense of Behrend-Fantechi \cite{BF}.\footnote{Virasoro constraints are also expected for moduli spaces of sheaves on Calabi-Yau 4-folds, admitting a virtual fundamental class in the sense of Oh-Thomas \cite{OT}. However, we will not consider this case here and it will be the subject of future work.} 

The main objects in this paper are moduli spaces $M$ which parameterize semistable sheaves on $X$ and their cohomology ring. Throughout this Introduction and Section \ref{sec: virasoro constraints}, when we refer to a general moduli space $M$ we assume that:
\begin{enumerate}
\item $M$ is a projective scheme of finite type. 
\item There are no $\BC$-points of  $M$ corresponding to strictly semistable sheaves.
\item Deformation theory at $[G]\in M$ is given by
\begin{align*}
\textup{Tan}&=\Ext^1(G,G)\\
\textup{Obs}&=\Ext^2(G,G)\\
0&=\Ext^{>2}(G,G)\,.
\end{align*}
\end{enumerate}


In such conditions, $M$ admits a virtual fundamental class $[M]^\vir\in H_\bullet(M)$ by \cite{BF}. There are many examples of moduli of sheaves on curves, surfaces and Fano 3-folds where all the assumptions are satisfied. To facilitate the exposition, for the introduction we will also assume that there exists a universal sheaf $\BG$ in $M\times X$; we will explain why this is not necessary in Section \ref{sec: invariantdescendents}. Note that a universal sheaf, when it exists, is non-unique: given a universal sheaf $\BG$ and a line bundle $L$ on $M$, the sheaf $\BG\otimes p^\ast L$ is also a universal sheaf (where $p\colon M\times X\to M$ is the projection) parametrizing the same objects.

Later in the paper we will also formulate Virasoro constraints for moduli spaces $M$ that have strictly semistable sheaves, by using Joyce's invariant class 
\[[M]^\inva \in H_\bullet(\CM_X^\rig)\,,\]
which is defined by a wall-crossing formula as we will overview in Section \ref{sec: joyceclasses}. This generalization will play an important role in the inductive argument used to prove Theorem \ref{thm: main}.
\begin{remark}
\source{virasoro-constraints-moduli-sheaves-vertex-algebras}
\notready
\label{rem:nocoarsedirect}
\source{virasoro-constraints-moduli-sheaves-vertex-algebras}
    The relation between the coarse moduli spaces $M$ and their invariant classes $[M]^{\inva}$ is not obvious in the presence of strictly semistables because there need not be a map $M\to \CM_X^{\rig}$. So saying that $M$ satisfies Virasoro constraints is an abuse of terminology that we will repeatedly make throughout this work. 
\end{remark}


Apart from the moduli of sheaves, we also study those of pairs. We fix a sheaf $V$ on $X$ and we let $P$ be a moduli space parametrizing a sheaf $F$ together with a map $V\to F$ (with some stability condition). We make the following assumptions: 
\begin{enumerate}
\item $P$ is a projective scheme of finite type. 
\item There are no $\BC$-points of $P$ corresponding to strictly semistable pairs.
\item There is a unique universal pair $q^*V\rightarrow\BF$ in $P\times X$ where $q\colon P\times X\to X$ is the projection. 
\item Deformation theory at $[V\to F]\in P$ is given by
\begin{align*}
\textup{Tan}&=\Ext^0([V\to F],F)\\
\textup{Obs}&=\Ext^1([V\to F],F)\\
0&=\Ext^{>1}([V\to F],F)\,.
\end{align*}
\end{enumerate}
In such conditions, $P$ admits a virtual fundamental class $[P]^\vir\in H_\bullet(P)$. Various moduli of pairs on curves and surfaces satisfy these assumptions; Quot schemes with at most one-dimensional quotients and moduli of Bradlow pairs.

There are two important differences between moduli spaces of sheaves and pairs.
\begin{enumerate}
\item The first is the difference in obstruction theory. It is apparent from comparing Example \ref{Ex: Tvir} and the definition of the Virasoro operators in Section \ref{sec: Virarosorops} that obstruction theory dictates their form.
\item  The second is the uniqueness or non-uniqueness of the universal object. This difference will play a crucial role in our treatment of the Virasoro constraints for moduli of sheaves and for moduli of pairs.
\end{enumerate}



\begin{remark}
\source{virasoro-constraints-moduli-sheaves-vertex-algebras}
\notready\label{rem: tracelessobs}
\source{virasoro-constraints-moduli-sheaves-vertex-algebras}
When we refer to moduli of sheaves we are mostly thinking about moduli of sheaves without fixed determinant. This is implicit in the obstruction theory above since when the determinant is fixed the deformation theory should instead use traceless Ext groups:
\[\Ext^i(G,G)_0=\ker\left(\Ext^i(G,G)\to H^i(\CO_X)\right).\]
We explain how to obtain a fixed determinant version of the Virasoro constraints in Section \ref{sec: fixeddet} when $h^{1,0}\neq 0$ but $h^{p,0}=0$ for $p>1$. Although a conjecture for Hilbert schemes of points on surfaces with possibly $p_g=h^{2,0}>0$  (which have traceless deformation theory) appears in \cite{moreira}, our approach in this paper is currently not suitable to understand it. We hope to pursue this direction in the future. 
\end{remark}


\begin{remark}
\source{virasoro-constraints-moduli-sheaves-vertex-algebras}
\notready\label{rem: stablepairsPT}
\source{virasoro-constraints-moduli-sheaves-vertex-algebras}
Virasoro constraints that we study for moduli of sheaves naturally generalize to moduli of objects in a derived category $D^b(X)$. Indeed, moduli spaces of stable pairs on a 3-fold $X$ (with $H^i(\CO_X)=0$ for $i>0$) in the sense of Pandharipande-Thomas \cite{PT} are instances of such. We emphasize here that stable pairs on 3-folds are subject to Virasoro constraints of sheaf type rather than pair type, despite their name. This is because virtual classes are constructed using the obstruction theory governed by $\Ext^i(I^\bullet,I^\bullet)$ where $I^\bullet=[\CO_X\to F]\in D^b(X)$. 

\end{remark}


\subsection{Universal sheaves and descendents}
\label{sec: descendents}
\source{virasoro-constraints-moduli-sheaves-vertex-algebras}
Descendents on $M$ are defined using a slant product construction with a universal sheaf $\BG$ and the maps

\begin{center}\begin{tikzcd}
& M\times X\arrow[ld, swap, "p"]\arrow[rd, "q"] &\\
M & & X\,.
\end{tikzcd}
\end{center}

\begin{definition}
\source{virasoro-constraints-moduli-sheaves-vertex-algebras}
\notready\label{def: geometricrealization}
\source{virasoro-constraints-moduli-sheaves-vertex-algebras}
We let $\BD^X$ be the supercommutative algebra generated by symbols $\chh_i(\gamma)$ for $i\geq 0$, $\gamma\in H^\bullet(X)$ (see Definition \ref{def: descendentalgebra}). The geometric realization with respect to a universal sheaf $\BG$ in $M\times X$ is the algebra homomorphism 
\[\xi_\BG\colon \BD^X\to H^\bullet(M)\]
defined on generators $\chh_i(\gamma)$ with $\gamma\in H^{r,s}(X)$ by 
\[\xi_\BG\left(\ch_i^\HH(\gamma)\right)=p_\ast\left(\ch_{i+\dim(X)-r}(\BG)\,q^\ast \gamma\right).\]
\end{definition}
The shift in the index of the Chern character using the Hodge degree of $\gamma$ is non-standard, but useful for a cleaner formulation of the Virasoro operators. With this convention, we may think of $\xi_\BG(\ch_i^\HH(\gamma))$ as being in $H^{i, i-r+s}(M)$ (of course $M$ might be singular, so a Hodge decomposition may not exist). See also Remark \ref{rem: comparingnotation}. 

The main objects of study in this paper are descendent integrals, i.e., the enumerative invariants obtained by integrating descendents against the virtual fundamental classes
\[\int_{[M]^\vir}\xi_{\BG}(D)\,,\quad \int_{[P]^\vir}\xi_{\BF}(D)\quad \textnormal{for}\quad D\in \BD^X.\]
Note that the descendent invariants of $M$ depend in principle on the choice of universal sheaf $\BG$. For some $D$ in the descendent algebra, however, they do not depend on this choice; these $D$ form what we call the  \textit{weight 0 descendent} algebra $\BD^X_{\inv}$ (cf. Section~\ref{sec: invariantdescendents}). For $D\in \BD^X_{\inv}$ we will omit the geometric realization morphism and write
\[\int_{[M]^\vir}D=\int_{[M]^\vir}\xi_\BG(D)\]
for any universal sheaf $\BG$ since it does not depend on such choice.

The Virasoro constraints say that these numbers satisfy some explicit universal relations. These relations are stated using certain operators
\[\bL_{\inv}\colon \BD^X\to \BD^X_{\inv}\,,\quad \bL_k^V:\BD^X\to \BD^X\,,\quad k\geq -1\]
that we will introduce in Section \ref{sec: virasoro constraints}. 
\begin{conjecture}[Virasoro for sheaves]
\source{virasoro-constraints-moduli-sheaves-vertex-algebras}
\notready\label{conj: virasorosheaves}
\source{virasoro-constraints-moduli-sheaves-vertex-algebras}
Let $M$ be a moduli of sheaves as before. Then
\[\int_{[M]^\vir} \Linv(D)=0\quad \textup{ for any }D\in \BD^X.\]
\end{conjecture}

\begin{conjecture}[Virasoro for pairs]
\source{virasoro-constraints-moduli-sheaves-vertex-algebras}
\notready
\label{conj: virasoropairs}
\source{virasoro-constraints-moduli-sheaves-vertex-algebras}
Let $P$ be a moduli of pairs as before. Then
\[\int_{[P]^\vir} \xi_\BF\left(\bL_k^V(D)\right)=0\quad \textup{ for any }k\geq 0, \, D\in \BD^X.\]
\end{conjecture}

In Example \ref{ex: rank2} we illustrate very explicitly how the constraints look like in the case of rank 2 stable bundles over a curve. 

\begin{remark}
\source{virasoro-constraints-moduli-sheaves-vertex-algebras}
\notready
The previous Virasoro conjectures for sheaves in \cite{moop, moreira, bree} require a specific choice of a universal sheaf and $\bS_k$ operators.\footnote{The $\bS_k$ operators in \cite{moop, moreira, bree} do not satisfy Virasoro bracket relations, obscuring the meaning of Virasoro constraints.} Conjecture \ref{conj: virasorosheaves} improves the formulation by avoiding both of these, and  we prove that the two formulations are equivalent (see Proposition \ref{prop: normalizedvsinv}). Conjecture \ref{conj: virasoropairs} for pairs is new and we provide convincing evidences by proving it for various geometries in this paper. 
\end{remark}

\subsection{Joyce's vertex algebra}
\label{sec: introvertexalgebra}
\source{virasoro-constraints-moduli-sheaves-vertex-algebras}


D. Joyce recently introduced a vertex algebra and a closely related Lie algebra associated to the derived category $D^b(X)$ \cite{Jo17,GJT, Jo21}. Joyce proposes to use his Lie algebra to study wall-crossing formulae for moduli of sheaves (or, more generally, moduli of semistable objects in a $\BC$-linear abelian or triangulated category).

The vertex algebra is constructed using the homology of the (higher) moduli stack $\CM_X$ parametrizing objects in the triangulated category $D^b(X)$. He defines a vertex algebra structure on
\[V_\bullet=\widehat{H}_\bullet(\CM_X)\,,\]
where $\widehat{H}_\bullet$ is meant to denote an appropriate shift in the grading of the homology.
The two most important ingredients for a vertex algebra are a translation operator $T$ and a state-field correspondence $Y(-,z)$; we will recall the definition of vertex algebras in Section \ref{sec: voadefinition}. In our setting, the translation operator is obtained from the $B\BG_m$--action on $\CM_X$; the state-field correspondence $Y(-,z)$ is defined in terms of the map 
$\Sigma:\M_X\times \M_X\to \M_X$ induced by taking direct sums and a perfect complex $\Theta$ on $\M_X\times \M_X$. These arise as a consequence of the master space localization technique that is commonly used for the proof of wall-crossing formulae (for instance in Mochizuki's work \cite{mochizuki}); the complex $\Theta$ is closely related to the virtual normal bundle appearing in the localization formula and thus to the obstruction theory of $\CM_X$. Remark \ref{rem: obstructionvsvirasoro} uses this observation to explain the relation of Virasoro constraints to the obstruction theory which we eluded to earlier on. 

Associated to the vertex algebra $V_\bullet$ is the Lie algebra obtained as the quotient by the translation operator:
\[\widecheck V_{\bullet}= V_{\bullet+2}/T V_{\bullet}\,.\]
The Lie bracket on $\widecheck V_\bullet$ is a shadow of the vertex algebra structure on $V_\bullet$ and is obtained by a well-known construction due to Borcherds \cite{Borcherds}.
Alternatively, it can be constructed as the homology $H_\bullet(\CM^\rig)$ of the rigidification $\CM_X^\rig=\CM_X \mkern-7mu\fatslash B\BG_m$; the two definitions agree when restricted to complexes with non-trivial numerical class, see Lemma \ref{Lem: DXinv}. 

The Lie algebra $\widecheck V_\bullet$ is a natural place where we can compare virtual fundamental classes of moduli spaces of sheaves; given a moduli space $M$ of semistable sheaves (or more generally of objects in $D^b(X)$) containing no strictly semistable sheaves, there is an open embedding $M\hookrightarrow \CM_X^\rig$. If $M$ admits a virtual fundamental class, we may push it forward along this embedding to obtain a class
\[[M]^\vir\in \widecheck{H}_\bullet(\CM_X^\rig)=\widecheck V_\bullet\]
where the first appearance of $\widecheck{(-)}$ represents an appropriate degree shift. If we fix a choice of a universal sheaf $\BG$ in $M\times X$, by the universal property of $\CM_X$ we get a map $f_\BG\colon M\to \CM_X$ lifting $M\hookrightarrow \CM_X^\rig$, and thus a natural lift of the virtual fundamental class to the vertex algebra
\[[M]^{\vir}_{\BG}\coloneqq (f_\BG)_\ast [M]^\vir\in V_\bullet\,.\]
Crucially, Joyce defines more general classes
\[[M]^\inva\in \widecheck V_\bullet\] 
even when strictly semistable sheaves exist; when $[M]^\vir$ is defined, both classes agree. 

The classes $[M]^\vir\in \widecheck V_\bullet$ or $[M]^{\vir}_{\BG}\in V_\bullet$ contain essentially the information of the (invariant) descendent integrals on $M$. This is made precise by J. Gross' \cite{Gr19} explicit description of $V_\bullet$, which we recall in Section \ref{sec: grossiso}. The cohomologies $H^\bullet(\CM_X)$ and $H^\bullet(\CM_X^\rig)$ are closely related to the algebras of descendents $\BD^X$ and $\BD^X_{\inv}$, respectively; see Lemmas \ref{Lem: isomorphism} and \ref{Lem: DXinv} for the precise statements. The pairing between cohomology and homology then recovers the descendent integrals
$$H^\bullet(\CM_X)\otimes H_\bullet(\CM_X)\to \BC
\,,\quad (D,[M]^{\vir}_{\BG})\mapsto \int_{[M]^\vir}\xi_{\BG}(D)\,,
$$
and 
$$H^\bullet(\CM_X^\rig)\otimes H_\bullet(\CM_X^\rig)\to \BC\,,\quad
(D,[M]^\vir)\mapsto \int_{[M]^\vir}D \,,
$$
where the second integral is independent of the choice of $\BG$.
%\begin{align*}
%H^\bullet(\CM_X)\otimes H_\bullet(\CM_X)&\to \BC\\
%(D,[M]^{\vir}_{\BG})&\mapsto %\int_{[M]^\vir}\xi_{\BG}(D)\,,
%\end{align*}
%and
%\begin{align*}
%H^\bullet(\CM_X^\rig)\otimes %H_\bullet(\CM_X^\rig)&\to \BC\\
%(D,[M]^\vir)&\mapsto \int_{[M]^\vir}D \,,
%\end{align*}




\subsection{Conformal element and Virasoro constraints}
\label{sec: introconformal}
\source{virasoro-constraints-moduli-sheaves-vertex-algebras}

A vertex operator algebra is a vertex algebra $V_\bullet$ equipped with a conformal element $\omega\in V_4$. The main property of a conformal element (see Section \ref{sec: voadefinition} for a precise definition) is that the operators $\{L_n\}_{n\in \BZ}$ on $V_\bullet$ induced from $\omega$ via the state-field correspondence satisfy the Virasoro bracket
\[\big[L_n,L_m\big] = (n-m)L_{m+n} + \frac{n^3-n}{12}\delta_{n+m,0}\cdot C\]
for some constant $C\in \BC$ called the central charge of $(V_\bullet, \omega)$. One of the goals of this paper is to explain the Virasoro operators in the descendent algebra previously studied in terms of a conformal element $\omega$ in Joyce's vertex algebra (or some slight variation, namely the pair vertex algebra). Due to the mysterious role that the Hodge degrees play in the Virasoro operators in \cite{moreira}, we do not know how to do so in complete generality, but only under the following assumption:


\begin{assumption}
\source{virasoro-constraints-moduli-sheaves-vertex-algebras}
\notready
\label{Ass: vanishing}
\source{virasoro-constraints-moduli-sheaves-vertex-algebras}
We assume that the Hodge cohomology groups $H^{p,q}(X)$ vanish whenever $|p-q|>1$. 
\end{assumption}
This assumption is satisfied for curves, surfaces with $p_g=0$ and Fano 3-folds, hence covering the majority of the target varieties in Donaldson-Thomas theory.

The result of J. Gross \cite{Gr19}\footnote{Some corrections to the statements in \cite{Gr19} were necessary for purposes of accuracy. We rederive all of the necessary results about the vertex algebra structure in Theorem \ref{thm: gross} to make sure that they are true.} shows that, under certain assumption (satisfied for curves, surfaces and rational 3-folds), Joyce's vertex algebra $V_\bullet$ is naturally isomorphic to a lattice vertex algebra from $(K^\bullet(X), K_\sst^0(X), \chi_\sym)$; here
\[K^\bullet(X)=K^0(X)\oplus K^1(X)\cong H^\bullet(X)\] is the topological $K$-theory of $X$ with $\BC$-coefficient, $K_\sst^0(X)$ is the semi-topological $K$-theory\footnote{The zeroth semi-topological $K$-theory $K^0_\sst(X)$ is defined as a Grothendieck's group of vector bundles modulo algebraic equivalence. By \cite[Theorem 4.21]{B16} and \cite[Theorem 2.3]{AH18}, we have $K^0_\sst(X)\simeq \pi_0(\mathcal{M}_X)$.} with $\BZ$-coefficient and $\chi_\sym$ is the symmetric pairing on $K^\bullet(X)$
\[\chi_\sym(v,w)=\int_{X}\ch(v^\vee)\ch(w)\td(X)+\int_{X}\ch(w^\vee)\ch(v)\td(X)\,.\]
The construction of a vertex algebra from such data is recalled and summarized in Theorem \ref{thm: voaconstruction} and follows Kac \cite{Ka98}; it uses Kac's bosonic vertex algebra construction in the even part $K^0(X)$ and the anti-fermionic vertex algebra construction in the odd part $K^1(X)$. 

Kac's construction produces a conformal element when the pairing $\chi_\sym$ is non-degenerate; unfortunately, due to the symmetrization this is not often the case. It turns out that this issue can be overcome by using the larger vertex algebra~$V_\bullet^{\pa}$. The vector space underlying $V_\bullet^{\pa}$ is the homology of the stack of pairs $\CP_X\simeq \CM_X\times \CM_X$:
\[V_\bullet^{\pa}=\widehat{H}_\bullet(\CP_X)\,.\]
The construction of the conformal element requires a choice of an isotropic decomposition of the fermionic part (see Theorem \ref{thm: voaconstruction} or \cite[Section 3.6]{Ka98}). This decomposition is where the Hodge degrees and Assumption \ref{Ass: vanishing} come into play, because $K^1(X)$ splits into the isotropic subspaces
$$K^1(X)=K^{\bullet,\bullet+1}\oplus K^{\bullet+1,\bullet}
$$
which via the Chern character isomorphism correspond to 
$$K^{\bullet,\bullet+1}\cong \bigoplus_{p\geq 0}H^{p, p+1}(X)\quad \textup{and}\quad K^{\bullet+1,\bullet}\cong \bigoplus_{p\geq 0}H^{p+1, p}(X)\,.
$$
%\[K^1(X)\cong H^{\textup{odd}}(X)=H^{\bullet, \bullet+1}(X)\oplus H^{\bullet+1, \bullet}(X)\,,\]
%where 
%\[H^{\bullet, \bullet+1}(X)=\bigoplus_{p\geq 0}H^{p, p+1}(X)\quad \textup{and}\quad  H^{\bullet+1, \bullet}(X)=\bigoplus_{p\geq 0}H^{p+1, p}(X)\,.\]
It is for the construction of such conformal element that we use the vertex algebra over the complex numbers while the result of J. Gross \cite{Gr19} works over any field containing rational numbers, as it relies on the Hodge decomposition. We prove that the Virasoro operators induced by this conformal element $\omega$ are dual to the pair Virasoro operators defined in the algebra of descendents, see Section \ref{sec: virasorocomparison}.


One remarkable aspect of Theorem B is that, while the operators $\bL_n^{\pa}$ on the descendent algebra were previously only defined for $n\geq -1$, a conformal element provides fields $L_n^{\pa}$ for every $n\in \BZ$ and thus a complete representation of the Virasoro algebra. In particular, this representation now has a non-trivial central charge $2\chi(X)$ which the positive branch $\{L_n^{\pa}\}_{n\geq -1}$ does not detect. We note that the factor of 2 appears due to working with the pair vertex algebra $V_\bullet^{\pa}$; if the pairing $\chi_\sym$ were non-degenerate we would get a conformal element in $V_\bullet$ with central charge $\chi(X)$. Remarkably, the Virasoro operators on the Gromov-Witten theory of $X$ (at least if $\dim(X)$ is even) are also known to admit an extension to a full representation of the Virasoro algebra; the central charge in the Gromov-Witten case is $\chi(X)$ \cite[Section 2.10]{getzler}. 

This description of the Virasoro operators provides a beautiful formulation of the Virasoro constraints for sheaves and for pairs in terms of well-known notions in the theory of vertex operator algebras, namely subspace of primary states (also known as physical states):
$$\widecheck{P}_0\subset \widecheck{V}_\bullet,\quad P_0^{\pa} \subset V_\bullet^{\pa}.$$

\begin{theorem}[Section \ref{subsec: virasoroprimary}]
\source{virasoro-constraints-moduli-sheaves-vertex-algebras}
\notready\label{thm: virasorophysical}
\source{virasoro-constraints-moduli-sheaves-vertex-algebras}
Assume $X$ is in class D (see Remark \ref{rem: classD}) and satisfies Assumption \ref{Ass: vanishing}. Then 
\begin{enumerate}
\item Conjecture \ref{conj: virasorosheaves} is equivalent to the class 
\[
[M]^\vir\in \widecheck{V}_\bullet \subseteq \widecheck{V}^{\pa}_\bullet
\]
being a primary state in $\widecheck{P}_0$ (cf. Definition \ref{Def: physicalstates}). 
\item Conjecture \ref{conj: virasoropairs} is equivalent to the class 
\[
[P]_{(q^\ast V,\BF)}^\vir\in V^{\pa}_\bullet
\]
being a primary state in $P_0^{\pa}$ 
(cf. Definition \ref{Def: physicalstates}). Here, the class $[P]_{{(q^\ast V,\BF)}}^\vir$ denotes the lift of $[P]^\vir\in \widecheck{V}^{\pa}_\bullet$ to $V^{\pa}_\bullet$ induced by the universal pair $q^\ast V\to\BF$.
\end{enumerate}
\end{theorem}
While the part (1) in the above theorem was stated for moduli spaces of sheaves satisfying the assumptions (1), (2) and (3) from Section \ref{sec: modulisheavespairs}, the condition of being a primary state makes sense also for the invariant classes $[M]^{\inva}$ without the assumption (2). This motivates the following generalization of Conjecture \ref{conj: virasorosheaves}. 
\begin{conjecture}
\source{virasoro-constraints-moduli-sheaves-vertex-algebras}
\notready
    \label{conj:actualinvvir}Under the assumptions of Theorem \ref{thm: virasorophysical}, the moduli spaces $M$, possibly containing strictly semistable sheaves, satisfy Virasoro constraints in the sense that
\source{virasoro-constraints-moduli-sheaves-vertex-algebras}
  %Joyce's invariant classes $[M]^{\inva}\in \widecheck{V}_{\bullet}\subset \widecheck{V}^{\pa}_{\bullet}$ are said to satisfy Virasoro constraints when they are primary with respect to the conformal element $\omega\in \widecheck{V}^{\pa}_{\bullet}$ discussed above, i.e. 
  $$
  [M]^{\inva}\in \widecheck{P}_0\subset  \widecheck{V}^{\pa}_{0}\,.
  $$
\end{conjecture}

The proof of Theorem \ref{thm: virasorophysical} is given in Section \ref{subsec: virasoroprimary}. We prove in Proposition \ref{prop: wallcrossingcompatibility1}, \ref{prop: wallcrossingcompatibility2} that the space of physical states interact nicely with Lie bracket operations: $\widecheck{P}_0\subset \widecheck{V}_\bullet
$ is a Lie subalgebra and $P_0^{\pa} \subset V_\bullet^{\pa}$ is a Lie submodule over $\widecheck{P}_0$. Since wall-crossing formulae in \cite{Jo21} are always written using the Lie bracket, these Lie algebraic statements prove a compatibility between wall-crossing and the Virasoro constraints. 


\subsection{Proof of Theorem \ref{thm: main} and other results}
\label{sec: results}
\source{virasoro-constraints-moduli-sheaves-vertex-algebras}

The main result of the paper is Theorem \ref{thm: main}, i.e., a proof of Conjecture \ref{conj: virasorosheaves} for semistable sheaves on curves and surfaces with $h^{0,1}=h^{0,2}=0$. We also denote the three cases (1), (2)  and (3) in Theorem \ref{thm: main} by $(m,d) = (1,1), (2,2)$ and $(2,1)$ where $m$ denotes the dimension of the underlying variety and $d$ the dimension of the support of sheaves we consider. 

The main ingredient in the proof is an inductive rank reduction argument via wall-crossing. This is the content of Section \ref{sec: rankreduction}. In each of the 3 cases, we consider the moduli spaces of Bradlow pairs $P^t_\alpha$ which depend on a stability parameter $t>0$. Assuming that $M_\alpha$ contains no strictly semistable sheaves, when $0<t\ll 1$ is small there is a map $P^{0+}_\alpha\coloneqq  P^{t\ll 1}_\alpha\to M_\alpha$ which is a (virtual) projective bundle. This geometric description is equivalent to wall-crossing at the Joyce-Song wall:
$$[P^{0+}_{\alpha}]^\vir = \big[[M_{\alpha}]^\vir,e^{(1,0)}\big]\,.$$
On the other hand, for large $t\gg 1$ the moduli spaces $P^{\infty}_\alpha\coloneqq  P^{t\gg 1}_\alpha$ are easier to understand (and sometimes empty). We prove not only that $M_\alpha$ satisfy the sheaf Virasoro constraints (i.e., Theorem  \ref{thm: main}) but also that the moduli spaces of Bradlow pairs $P_\alpha^t$ satisfy the pair analogue of the constraints:
\begin{theorem}[Theorem \ref{thm: rankreduction}, Section \ref{sec: limits}, Section \ref{sec:lowrank}]
\source{virasoro-constraints-moduli-sheaves-vertex-algebras}
\notready\label{thm: virasorobradlowpairs}
\source{virasoro-constraints-moduli-sheaves-vertex-algebras}
The moduli spaces of Bradlow pairs $P_\alpha^t$ (see Definition \ref{def: bradlow}) satisfy the pair Virasoro constraints (Conjecture \ref{conj: pairvirasoro}) for every $t>0$ in the 3 settings of Theorem \ref{thm: main}, i.e., $(m,d)=(2,2), (1,1)$ and $(m,d)=(2,1)$ provided Assumption \ref{ass:WCpair} holds.
\end{theorem}

To prove Theorems \ref{thm: main} and \ref{thm: virasorobradlowpairs} we need the following steps:
\begin{enumerate}
\item[(i)] We prove in Section \ref{sec:lowrank} that $P^{\infty}_\alpha$ satisfies Conjecture \ref{conj: virasoropairs}. In case $(m,d)=(1,1)$, we only need to prove the statement for symmetric powers of curves which we do by a direct computation in Proposition \ref{prop: symmetricpowers}. Cases $(m,d)=(2,1)$ and $(2,2)$ for slope stability can be reduced to Hilbert schemes of points, where they were shown to hold in \cite{moop, moreira}.
\item[(ii)] We use the wall-crossing formula \eqref{eq: rankreductionwallcrossing} between $P^{\infty}_\alpha$ and $P^{0+}_\alpha$ to show that $P^{0+}_\alpha$ satisfies Virasoro constraints for pairs as well. By induction on $\rk(\alpha)$, we know the Virasoro constraints on the wall-crossing terms. We then rely on the compatibility between wall-crossing and Virasoro constraints (Propositions \ref{prop: wallcrossingcompatibility1} and \ref{prop: wallcrossingcompatibility2}).
\item[(iii)]  Finally, we use a projective bundle compatibility for  $P^{0+}_\alpha\to M_\alpha$ proved in Theorem \ref{lem: projectivebundle} to show that the pair Virasoro constraints on $P^{0+}_\alpha$ imply the sheaf Virasoro constraints on $M_\alpha$.  Wall-crossing from slope stability to Gieseker stability in Corollary \ref{cor:slopetogies} concludes the proof of Theorem \ref{thm: main}.
\end{enumerate}
A crucial point in the argument is that we must include moduli spaces $M_\alpha$ admitting strictly semistable sheaves in the induction since they unavoidably appear as wall-crossing terms. That is, we must prove that $[M]^\inva$ is in the Lie algebra of primary states. Because of that, in step (iii) we do not exactly have a projective bundle. However, by the very definition of the invariant classes $[M]^\inva$, what we have to prove is essentially the same as in the projective bundle case. We do this in Theorem~\ref{lem: projectivebundle}.

In the appendix we explain Joyce-Song wall-crossing, using results of the first author in \cite{Bo21.5}, which provides an alternative for some of the arguments in Section \ref{sec:lowrank}. In particular, we prove that the pair Virasoro constraints hold for punctual Quot schemes.\footnote{It would be interesting to see the implication of Virasoro constraints to the generating series of descendent invariants of Quot schemes studied in \cite{JOP}.}

\begin{theorem}[=Theorem \ref{prop: quot}]
\source{virasoro-constraints-moduli-sheaves-vertex-algebras}
\notready\label{thm: punctualquot}
\source{virasoro-constraints-moduli-sheaves-vertex-algebras}
Let $X$ be a curve or a surface and let $V$ be a torsion-free sheaf on $X$. Then the punctual Quot scheme $\textup{Quot}_X(V, n)$ satisfies the pair Virasoro constraints (Conjecture \ref{conj: pairvirasoro}).
\end{theorem}


\subsection{Notation and conventions}
Except the semi-topological $K$-group $K_\sst^0(X)$ which we consider over $\BZ$-coefficients, all cohomology and $K$-theory groups are assumed to have coefficients in $\BC$ unless stated otherwise. We write 
\[K^\bullet(X)=K^0(X)\oplus K^1(X)\]
for the topological $K$-theory and we denote by
\[\ch\colon K^\bullet(X)\to H^\bullet(X)\]
the isomorphism between topological $K$-theory and cohomology. The total cohomology of a topological space is always understood to be the direct product of the cohomology groups in each degree, while homology is the direct sum. We will use the cap product with the cohomology acting on the left. This is the convention followed in \cite{dold}; it differs from the more usual convention with cohomology on the right by a sign, i.e., $\gamma \cap u=(-1)^{|\gamma||u|}u\cap \gamma$ for $\gamma\in H^\bullet(X)$, $u\in H_\bullet(X)$. 

By an obstruction theory on a scheme $M$ we mean a perfect complex $\BE$ of non-negative degree together with a map $\phi:\BE^\vee\to \BL_M$ such that $h^0(\phi)$ and $h^{-1}(\phi)$ are an isomorphism and a surjection, respectively. The complex $\BE$, or its $K$-theory class, is also called the virtual tangent bundle of $M$ (note that by \cite{siebert} the virtual fundamental class only depends on the $K$-theory class $\BE$, and in particular does not depend on the map to $\BL_M$). We will often abuse notation and just call $\BE$ (or its $K$-theory class) the obstruction theory, leaving the map to $\BL_M$ implicit. The construction of the morphism to $\BL_M$ in all the cases considered is standard, see \cite{huythomas} for sheaves/complexes and e.g. \cite[Section 5.3]{mochizuki} for pairs.


\begin{center}
    \begin{tabular}{p{4cm} p{10cm}}
    $\alpha,\beta$& Semi-topological $K$-theory classes in $K_\sst^0(X)$.\\
    $\gamma, \delta$&  Cohomology classes on $X$. \\
    $v,w$& Elements of $K^\bullet(X)$.\\
    $\deg(-)$ & Degree for any graded vector space. 
    \\
    $|-|$& Supergrading taking value in $\{0,1\}. $   \\
    $\chh_i(\gamma)$ & The holomorphic descendent in degree $2i-p+q$ depending on the Hodge degree of $\gamma\in H^{p,q}(X)$.\\
    $\ch_i(\gamma)$& The topological descendent in degree $2i-|\gamma|$. \\
     $L_n, T_n, R_n,$& Virasoro operators on homology and vertex algebra.\\
    $\mathsf{L}_n, \mathsf{T}_n,\mathsf{R}_n,$& Dual operator notation on cohomology and descendent algebra.\\
    $\mathsf{L}_{\inv}$& weight 0 Virasoro operator on descendent algebra.
        \end{tabular}
\end{center}

\subsection{Future directions} 
There are several open directions regarding the Virasoro constraints for sheaves. The first obvious direction is to try to improve Theorem \ref{thm: main} by removing the assumptions $h^{0,1}=h^{0,2}=0$. The arguments in this paper show that we can get the constraints for $h^{0,1}>0$ as long as we can prove them for the moduli of rank 1 sheaves (isomorphic to the Hilbert scheme of points times the Jacobian). Finding an argument that works in general for the Hilbert scheme of points and does not go through Gromov-Witten theory would be highly desirable. Removing the assumption $h^{0,2}=0$ requires a better understanding of the constraints in the setting of reduced virtual fundamental classes for fixed determinant theory (see Remark \ref{rem: fixed det}).

Sheaf-theoretic Virasoro constraints of Fano 3-folds are of particular interest because they are related to the original Virasoro constraints in Gromov-Witten theory. Since wall-crossing compatibility of Virasoro constraints also holds in this case, it would be interesting to develop wall-crossing techniques for Fano 3-folds that can be applied to Virasoro constraints.

\subsection{Acknowledgement} 
The authors would like to thank Y. Bae, D. van Bree, J. Gross, A. Henriques, D. Joyce, M. Kool, H. Liu, A. Mellit,  G. Oberdieck, A. Oblomkov, R. Pandharipande and W. Pi for many conversations about the Virasoro constraints on moduli spaces of sheaves.

The first and third authors were supported by
ERC-2017-AdG-786580-MACI. The second author is supported by the grant SNF-200020-182181. The project received funding from the European Research Council (ERC) under the European Union Horizon 2020 research and innovation programme (grant agreement 786580).
